\begin{tikzpicture}[scale=0.98, every node/.style={font=\small}]
  \draw[fill=softblue,draw=linegray] (0,0) rectangle (6.3,4.5);

  \coordinate (camera) at (5.7,4.1);
  \coordinate (impact) at (2.61,1.56);
  \coordinate (obstruction) at (3.67,2.44);

  \begin{scope}[rotate around={20:(3.25,1.7)}]
    \draw[fill=white,draw=deepblue,thick] (2.6,1.05) rectangle (3.9,2.35);
  \end{scope}

  % Le rayon laser s'arrête exactement sur la face visible du carré.
  \draw[heigred,very thick] (0.7,3.8) -- (impact);
  \draw[heigred,fill=heigred] (impact) circle (0.06);

  % La ligne pointillée est l'axe caméra -> point laser. La vision réelle
  % s'arrête avant ce point, sur la face qui masque la ligne laser.
  \draw[deepblue,very thick,dashed] (camera) -- (impact);
  \draw[deepblue,thin] (camera) -- (obstruction);
  \draw[deepblue,fill=deepblue] (obstruction) circle (0.035);

  \node[heigred,align=center] at (1.2,4.1) {laser};
  \node[deepblue,align=center] at (5.45,4.1) {caméra};
  \node[align=center] at (3.2,0.55) {face de l'objet};
  \node[align=center,text width=2.2cm] at (4.85,1.35) {zone masquée depuis la caméra};
\end{tikzpicture}
